% ---------- Preamble: shared setup, diagram and evidence macros ----------
% Policy: all body text, headings, captions, tables and callouts are BLACK.
% Colour appears ONLY inside TikZ diagrams.
\usepackage[T1]{fontenc}
\usepackage[utf8]{inputenc}
\usepackage[a4paper,margin=1in]{geometry}
\usepackage{graphicx}
\usepackage{xcolor}
\usepackage{booktabs}
\usepackage{colortbl}
\usepackage{array}
\usepackage{tabularx}
\usepackage{enumitem}
\usepackage{longtable}
\usepackage{float}
\usepackage{caption}
\usepackage{amssymb}
\usepackage{fancyhdr}
\usepackage{titlesec}
\usepackage{listings}
\usepackage{tikz}
\usetikzlibrary{positioning,arrows.meta,calc,fit,backgrounds,shapes.geometric}
\usepackage[hidelinks]{hyperref}

% ---------- Diagram-only palette (never used for text) ----------
\definecolor{accent}{HTML}{2C6E91}
\definecolor{accentlt}{HTML}{E6EEF3}
\definecolor{good}{HTML}{2E7D4F}
\definecolor{goodlt}{HTML}{E4F1E9}
\definecolor{bad}{HTML}{B23A3A}
\definecolor{badlt}{HTML}{F6E5E4}
\definecolor{net1}{HTML}{E7F0F7}
\definecolor{net2}{HTML}{FBEFE3}
\definecolor{net3}{HTML}{EAF2EA}
\definecolor{dslate}{HTML}{45566B}
% grayscale tones for placeholders and code frames (neutral, not colour)
\definecolor{boxframe}{gray}{0.0}
\definecolor{boxfill}{gray}{0.96}
\definecolor{boxtext}{gray}{0.25}
\definecolor{coderule}{gray}{0.55}

% ---------- Section styling: black ----------
\titleformat{\section}{\color{black}\Large\bfseries}{\thesection}{0.6em}{}
\titleformat{\subsection}{\color{black}\large\bfseries}{\thesubsection}{0.6em}{}
\titleformat{\subsubsection}{\color{black}\normalsize\bfseries}{\thesubsubsection}{0.6em}{}
\titlespacing*{\section}{0pt}{1.4em}{0.7em}

% ---------- Header / footer: black ----------
\pagestyle{fancy}
\fancyhf{}
\renewcommand{\headrulewidth}{0.4pt}
\renewcommand{\footrulewidth}{0pt}
\fancyhead[L]{\small CS6008 Network Security}
\fancyhead[R]{\small Assignment 2: DNS Security}
\fancyfoot[C]{\small\thepage}

% ---------- Listings: black text, white ground, plain frame ----------
\lstdefinestyle{term}{
  basicstyle=\ttfamily\footnotesize\color{black},
  backgroundcolor=\color{white},
  breaklines=true,
  breakatwhitespace=false,
  columns=fullflexible,
  keepspaces=true,
  frame=single,
  rulecolor=\color{coderule},
  framesep=6pt,
  xleftmargin=10pt,
  xrightmargin=4pt,
  aboveskip=10pt,
  belowskip=8pt,
  showstringspaces=false
}
\lstset{style=term}

% ---------- Evidence figures ----------
% Each figure embeds its screenshot AUTOMATICALLY, found by the figure's label:
% a figure with \label{fig:p2-attack} looks for an image called p2-attack.
% It searches, in order:  images/p2-attack.png, images/p2-attack.PNG,
% images/p2-attack.jpg, images/p2-attack.jpeg, images/p2-attack.pdf, and the
% same five names at the project root.  As soon as one of those files exists
% the picture appears; until then a gray "reserved space" box is shown that
% prints the exact filename it is waiting for.  So the ONLY thing you have to do
% is name each screenshot after its figure (see the list in SCREENSHOTS.md) and
% drop them in an images/ folder.  Nothing in report.tex needs editing.
\graphicspath{{images/}{./}}

% strip the leading "fig:" from a label so we can build the image filename
\def\stripfigcolon fig:#1\stripfigend{#1}
\newcommand{\figimgname}[1]{\expandafter\stripfigcolon#1\stripfigend}

% try to include <base> with any known folder+extension; run #2 if none found
\newcommand{\tryshot}[2]{% #1 = basename (e.g. p2-attack), #2 = fallback code
  \IfFileExists{images/#1.png}{\includegraphics[width=0.92\linewidth]{images/#1.png}}{%
  \IfFileExists{images/#1.PNG}{\includegraphics[width=0.92\linewidth]{images/#1.PNG}}{%
  \IfFileExists{images/#1.jpg}{\includegraphics[width=0.92\linewidth]{images/#1.jpg}}{%
  \IfFileExists{images/#1.jpeg}{\includegraphics[width=0.92\linewidth]{images/#1.jpeg}}{%
  \IfFileExists{images/#1.pdf}{\includegraphics[width=0.92\linewidth]{images/#1.pdf}}{%
  \IfFileExists{#1.png}{\includegraphics[width=0.92\linewidth]{#1.png}}{%
  \IfFileExists{#1.jpg}{\includegraphics[width=0.92\linewidth]{#1.jpg}}{%
  #2}}}}}}}%
}

% #1 = caption, #2 = what the screenshot should show, #3 = suggested file,
% #4 = label. #5 = minipage height for the fallback box.
\newcommand{\shotboxinner}[5]{%
\begin{figure}[H]
\centering
\edef\figbase{\figimgname{#4}}%
\tryshot{\figbase}{%
  \setlength{\fboxsep}{0pt}%
  \fcolorbox{boxframe}{boxfill}{%
  \begin{minipage}[c][#5][c]{0.90\linewidth}
  \centering
  {\small\itshape\color{boxtext} Space reserved for screenshot}\\[3pt]
  {\small\color{boxtext} #2}\\[5pt]
  {\footnotesize\color{boxtext} Save screenshot as \texttt{images/\figbase.png} \quad (source: \texttt{#3})}
  \end{minipage}}%
}%
\caption{#1}
\label{#4}
\end{figure}}

\newcommand{\shot}[4]{\shotboxinner{#1}{#2}{#3}{#4}{40mm}}
\newcommand{\shottall}[4]{\shotboxinner{#1}{#2}{#3}{#4}{58mm}}

% ---------- Callout: black framed box ----------
\newcommand{\takeaway}[1]{%
\par\smallskip\noindent
\setlength{\fboxsep}{7pt}%
\fbox{\parbox{\dimexpr\linewidth-2\fboxsep-2\fboxrule\relax}{\small #1}}%
\par\smallskip}

\newcolumntype{Y}{>{\raggedright\arraybackslash}X}
\captionsetup{font=small,labelfont=bf,textfont=normal}